%% kapa: exam=1 round=34 year=2023 period=2 form=A subject=accounting subject_ko=회계학 seq=38 number=78
%% source: Q-Net 감정평가사 기출문제 (cst00309) artlSeq=5211789 file=raw/1008/5211789/2023년도 제34회 감정평가사 1차시험 2교시.pdf
%% conversion: kordoc 4.24.51 -> tools/gen1.py (auto); review=proofread-sonnet; flags=none
\begin{question}{78}
(주)감평은 정상개별원가계산제도를 채택하고 있다. 제조간접원가는 직접노무원가의 40\%를 예정배부하고 있으며, 제조간접원가 배부차이는 전액 매출원가에서 조정하고 있다. (주)감평의 당기 재고자산 및 원가 관련 자료는 다음과 같다.
\par{\small\centering\begin{tabular}{|c|c|c|}\hline
구 분 & 기초잔액 & 기말잔액 \\ \hline
직접재료 & ￦3,200 & ￦6,200 \\ \hline
재 공 품 & 8,600 & 7,200 \\ \hline
제 \quad 품 & 6,000 & 8,000 \\ \hline
\multicolumn{3}{|l|}{ － 직접재료매입액 : ￦35,000} \\ \hline
\multicolumn{3}{|l|}{ － 기초원가(기본원가) : ￦56,000} \\ \hline
\end{tabular}\par}
(주)감평의 당기 제조간접원가 배부차이 조정 후 매출원가가 ￦67,700인 경우, 당기에 발생한 실제 제조간접원가는?
\begin{choices}
  \item ￦6,900
  \item ￦9,700
  \item ￦10,700
  \item ￦11,300
  \item ￦12,300
\end{choices}
\end{question}
\answer{5}
